% META: {"id":"1NSI-TC-EX-053","chapitre":"1NSI-TYPES-CONSTRUITS","type_objet":"exercice","capacites":["1NSI-TYPES-CONSTRUITS-C5"],"parcours":2,"competences":["analyser"],"duree_min":12,"mode_creation":"ex_nihilo","sources_inspiration":[],"status":"approved","origin":"needs_review","status_history":["needs_review","audited","accepted","approved"],"release_acceptance":"RELEASE_OWNER_BATCH_ACCEPTANCE","acceptance_closure_digest":"sha256:9b3ccf9a81c5520fbb7e03b7b2d3e7bf2057a02834b6d908bf3be5f49f3b553f","acceptance_receipt_digest":"sha256:4fad86f0282e0fa4e0419cca1ad6379ae7af5a280c04a4a90880f90a1c044242"}
% BEGIN-TRACE
% grille = [[1, 2], [3, 4]]
% copie = list(grille)
% copie[0][0] = 99
% print(grille[0][0])
% copie[0] = [10, 20]
% print(grille[0])
% EXPECTED
% 99
% [99, 2]
% END-TRACE
\begin{exercice}{1NSI-TC-EX-053}{2}{12}
On considère le programme suivant :
\begin{python}
grille = [[1, 2], [3, 4]]
copie = list(grille)
copie[0][0] = 99
print(grille[0][0])
copie[0] = [10, 20]
print(grille[0])
\end{python}
\begin{enumerate}
  \item Donner, sans exécuter, les deux affichages.
  \item Expliquer pourquoi \lstinline{grille[0][0]} vaut \lstinline{99} alors qu'on
        a modifié \lstinline{copie}.
  \item On parle de \textbf{copie superficielle}. Qu'est-ce que cela signifie ?
  \item On suppose désormais que les cellules sont des valeurs scalaires atomiques
        non mutables, sans conteneur imbriqué. Écrire une fonction
        \lstinline{copier_grille_deux_niveaux(grille)} qui réalise une
        \textbf{copie des deux premiers niveaux}. Sa garantie est limitée aux
        modifications de la structure externe, aux remplacements de lignes et aux
        réaffectations de cellules dans la copie, qui ne doivent pas modifier la
        grille d'origine.
  \item Hors de cette précondition, analyser le contre-exemple
        \lstinline{[[[1]]]}. Expliquer pourquoi la cellule-liste reste partagée
        entre l'original et la copie.
\end{enumerate}
% PYTHON-SOURCE: code/copier_grille_deux_niveaux.py
\end{exercice}
